\section{Related Work}
\label{sec:related}

\subsection{Honeypots and Cyber Deception}

Surveys of large language models in security \cite{wang2024} and of cyber
deception specifically \cite{lopez2024} frame the opportunity and its open
questions. Honeypots have been studied as instrumentation for two decades
\cite{spitzner2002}, and their engineering trade-offs---visibility,
fidelity, and the operational cost of interaction---remain open
\cite{sokol2017}. Adaptive deception work formalizes the
defender's decision to engage, treating attacker intent as something to
infer and modulate, through factored interactive POMDPs
\cite{shinde2020}, engagement--yield models
\cite{bera2024}, and behavioral partitioning under cognitive load
\cite{liska2024}. Deception has also been embedded in higher-fidelity
simulation and emulation targets \cite{kouremetis2024}. We adopt as our
starting point the gap that deception-framework surveys identify in deployed
systems: that they neither ground responses in live threat intelligence
nor differentiate behavior across attacker classes \cite{lopez2024}. Our
contribution is orthogonal to that taxonomy. We do not propose a new
deception policy; we constrain what a generative policy is allowed to serve.

\subsection{Generative Honeypots}

The closest prior art puts a language model inside the deception endpoint.
LLM-in-the-shell honeypots make the model answer as a live machine
\cite{valeros2024}; LLM agent honeypots instead monitor autonomous AI
attackers in the wild \cite{reworr2024}; and recent systems extend
generative deception to industrial protocols and physical-process emulation
\cite{vasilatos2025llmpot}, coordinated IoT deception
\cite{guan2025cpds}, and directory services
\cite{jimenez2025}. Analysis pipelines over collected honeypot and
incident data apply generative or learned models after the fact, for attack
classification and reporting \cite{lanka2024,simoni2024,blefari2025},
structured CTI mapping \cite{orbinato2022}, and attack investigation
\cite{alsaheel2021}. In all of these systems the model's output is the product, and the
product ships unchecked. None of them states a serving contract for that
output---no fixed fiction to be consistent with, no deterministic gate
that can refuse a draft, no vetted answer to fall back to when the
provider refuses or times out. To our knowledge, that contract is what
this paper adds, together with the ladder analysis that prices it; the
generative mechanisms themselves are borrowed.

\subsection{Runtime Enforcement and Red-Teaming of LLM Applications}

The verifier is the closest analogue in prior art, and the analogy is
deliberate. AgentSpec provides runtime enforcement for LLM agents
\cite{wang2026agentspec}; prompt-injection defences that constrain
decoding rather than trusting a probability appear in
\cite{chen2025struq,chen2025secalign}; the injection threat is
formalized in \cite{greshake2023} and measured by AgentDojo
\cite{debenedetti2024agentdojo}; tool-bearing agents add cross-tool
harvesting \cite{li2026dissonances}; model-red-teamed evaluation is
standard practice \cite{ganguli2022,perez2022,bai2022,christiano2017}. Our design follows
that literature's premise that assurance must live outside the model. What
differs is the asset being protected. A guardrail protects the assistant's
user from the model's output; this gate protects the attacker's belief from
the fiction's, so the failure of interest is a fluent, self-consistent,
on-script response that is nonetheless wrong for the deception state, and
the mitigation is a deterministic fallback to vetted content rather than a
refusal.

\subsection{Retrieval Grounding, and the Attacks We Do Not Yet Resist}

Retrieval grounding supplies the targeting rung \cite{lewis2020};
retrieval for threat intelligence is well populated---technique
extraction \cite{fayyazi2023,lekssays2025}, knowledge-graph construction
\cite{gao2025,jeon2024,belmahjoub2025}, index families \cite{malkov2018}
we price against sparse retrieval, response assistants
\cite{rajapaksha2024,paul2025,qiao2025,liu2025}, and the \attck{}
grounding question \cite{roy2023,alsada2024,noureddine2020}. The
security of that retrieval layer is an active literature: vulnerability
identification \cite{xue2024,liang2025}, RAG threat vectors and
mitigations \cite{ward2025}, knowledge corruption
\cite{zou2025poisonedrag,chaudhari2024}, retrieval jamming
\cite{shafran2024}, poisoning evaluation \cite{malladi2026ripeii},
poisoning-aware CTI ranking \cite{jia2025}, and attribution for poisoned
knowledge \cite{zhang2026ragorigin}. We adopt that literature's threat
model as out of scope for demonstrated results: this paper evaluates no
retrieval poisoning and no injection defence, and
Section~\ref{sec:threat-model} records that as a gap. The retrieval layer
sits inside the threat surface this literature has already mapped, so a
deployment that inherits it inherits those attacks, and the grounding
rung's measured benefit of $+0.154$ is a benefit under an undefended
retrieval layer. Closing
that is a prerequisite for treating the grounding result as a security
contribution rather than a capability one.

\subsection{Deception Quality Judged by Machines}

Because no human study is claimed, the relevant prior art is LLM-as-judge
methodology: MT-Bench and Chatbot Arena for judge agreement
\cite{zheng2024}, JudgeBench for judge reliability \cite{tan2025judgebench},
and optimization-based attacks against the judge itself
\cite{shi2024judgedeceiver}. The third of these is the reason our quality
section reports within-model consistency with an audited judge rather than
presenting a realism score as validated, and the reason our labeling lane
is disclosed as an AI panel with a deterministic-oracle ceiling
check. We also contribute the check itself: a blinded echo audit that
invalidated 97 of our own 218 raw scoring rows, offered as a reusable
instrument for any LLM-scored deception evaluation.

\subsection{Position}

We are not the first to build a language model into a honeypot, and we do
not claim a new deception policy, a new retrieval stack, or a new
enforcement primitive. The contribution is the composition: a fixed
deception state, a deterministic gate with authority to refuse, a vetted
floor beneath failure, and an explicit ladder analysis that prices each
intelligence mechanism against its cheapest certifiable competitor---then
removes the rungs that do not pay. The next section runs that pricing on
the machine itself. Because the close prior art evaluates
generation rather than serving, the battery of Section~\ref{sec:evaluation}
is also the first evaluation apparatus for the failure mode every
generative honeypot shares: the fluent, self-consistent, on-script draft
that is wrong for the deception state.